\section{The ReLexTail Order}
\label{sec:relextail}

Before defining the new rule, it is necessary to identify why the exact lexicographic minimax fails as a stable selection operator.

\subsection{Deficiencies of the exact lexicographic minimax}
The exact lexicographic minimax suffers from four fundamental limitations:
\begin{enumerate}
    \item \textbf{Infinite priority of the maximum:} If $D^\downarrow_1(x) < D^\downarrow_1(y)$, then the exact rule strictly prefers $x$, regardless of the magnitude of the difference $D^\downarrow_1(y)-D^\downarrow_1(x)$. This remains true even if the difference is arbitrarily small.
    \item \textbf{Instability near ties:} When the maximum disappointments are nearly equal, an arbitrarily small perturbation in the bounds may reverse the recommendation before the rest of the profile is even considered.
    \item \textbf{Poor compensation within the upper tail:} The exact rule cannot recognise that a slight deterioration in one maximum concern may accompany large improvements in the next several worst concerns.
    \item \textbf{Pairwise tolerance is nontransitive:} While an $\tau$-tolerance can be introduced, the resulting pairwise relation is nontransitive because approximate equality is not an equivalence relation, creating cycles.
\end{enumerate}

\subsection{Design Requirements}
The new order, ReLexTail (Resolution-Aware Lexicographic Tail-Regret Selection), satisfies the following requirements (Table~\ref{tab:design}):
\begin{table*}[t]
\centering
\caption{Design requirements of ReLexTail compared to previous formulations.}
\label{tab:design}
\begin{tabular}{@{}P{0.25\textwidth}P{0.18\textwidth}P{0.25\textwidth}P{0.20\textwidth}@{}}
\toprule
Requirement & Exact leximax & Pairwise tolerance & ReLexTail \\
\midrule
Complete & Yes & No/conditional & Yes \\
Transitive & Yes & No & Yes \\
Pareto-compatible & Yes & Conditional & Yes \\
Resolution-aware & No & Yes & Yes \\
Upper-tail-sensitive & Only after exact ties & Only after approx. ties & Yes \\
Uniform replication-invariant & Yes & Unclear & Yes \\
Auditable & Yes & Complicated by cycles & Yes \\
Interval-certifiable & Yes & Difficult & Yes \\
\bottomrule
\end{tabular}
\end{table*}

\subsection{Probe-disappointment profile and Upper-Tail Means}
Let $D(x) = (D_1(x), \ldots, D_K(x)) \in [0,1]^K$ denote the disappointment profile of an alternative $x$. We sort the components in decreasing order:
\[ D_{[1]}(x) \geq D_{[2]}(x) \geq \cdots \geq D_{[K]}(x) \]
The maximum disappointment is $M(x) = D_{[1]}(x)$.

For a tail proportion $\alpha \in (0,1]$, define the empirical upper-tail conditional value-at-risk (CVaR):
\[ T_\alpha(x) = \operatorname{CVaR}^{\uparrow}_{\alpha}(D(x)) = \min_{t\in\mathbb R} \left[ t + \frac{1}{\alpha K} \sum_{q=1}^{K} \max\{D_q(x)-t, 0\} \right]. \]
The probe coordinates are treated as equally weighted atoms of an empirical disappointment distribution. For finite computation with $h = \alpha K$, $k = \lfloor h \rfloor$, and $\theta = h - k$:
\[ T_\alpha(x) = \frac{\sum_{j=1}^{k} D_{[j]}(x) + \theta D_{[k+1]}(x)}{h}, \]
with the conventions that an empty sum is zero, if $\theta=0$ the $D_{[k+1]}$ term is omitted, and if $h<1$ then $k=0$, $\theta=h$, and $T_\alpha(x) = D_{[1]}(x) = M(x)$. Note also that $T_1(x) = \frac{1}{K}\sum_{q=1}^K D_q(x)$.

\subsection{Resolution Categories and ReLexTail Profile}
We define a fixed risk spectrum $\mathcal A = \{0.25, 0.50, 1.00\}$ corresponding to the severe-concern tail $T_{25}(x)$, the broad-concern tail $T_{50}(x)$, and the overall disappointment $T_{100}(x)$.

Given resolutions $\delta_M, \delta_{25}, \delta_{50}, \delta_{100} > 0$, we define the resolution categories:
\[ B_\delta(z) = \left\lceil \frac{z}{\delta} \right\rceil. \]
Under this convention, $B_\delta(0) = 0$, and for $j \ge 1$, $B_\delta(z) = j \iff (j-1)\delta < z \le j\delta$. Thus, the order is discontinuous at category boundaries. This intentional design distinguishes sensitivity within a category, sensitivity near a category boundary, and sensitivity of the exact-tail refinement.
Lower categories are preferred. We compute $C_M(x) = B_{\delta_M}(M(x))$, $C_{25}(x) = B_{\delta_{25}}(T_{25}(x))$, $C_{50}(x) = B_{\delta_{50}}(T_{50}(x))$, and $C_{100}(x) = B_{\delta_{100}}(T_{100}(x))$.

The ReLexTail profile is defined as:
\[ \Psi(x) = \Bigl(C_M(x), C_{25}(x), C_{50}(x), C_{100}(x), T_{25}(x), T_{50}(x), T_{100}(x), M(x), D_{[1]}(x), \ldots, D_{[K]}(x)\Bigr). \]
Note that $M(x) = D_{[1]}(x)$ is repeated. This repetition is harmless but redundant, present because $M$ has a specific position in the hierarchy while the complete sorted profile is used only for terminal refinement.
The alternatives are compared lexicographically:
\[ x \preceq_{\mathrm{ReLexTail}} y \iff \Psi(x) \leq_{\mathrm{lex}} \Psi(y). \]

This defines a clear hierarchy: catastrophic-resolution protection, followed by tail categories, exact tail refinements, exact maximum refinement, and finally complete deterministic refinement via the exact leximax. Every material category difference has priority over every within-category numerical difference. Furthermore, putting exact $M$ immediately after $C_M$ would recreate the exact-leximax priority inside every maximum category. ReLexTail therefore postpones exact maximum comparison until after broader tail comparisons.

\subsection{Resolution width versus pairwise difference tolerance}
It is essential to distinguish categorical resolution from pairwise tolerance. A pairwise tolerance requires that $|a-b| < \delta$ implies indifference, which famously violates transitivity. In ReLexTail, categories encode membership in fixed cells. The condition $|a-b| < \delta$ does not imply $B_\delta(a) = B_\delta(b)$, nor does $B_\delta(a) = B_\delta(b)$ imply that $|a-b|$ is arbitrarily small (except that $|a-b| \leq \delta$ under the ceiling convention). Furthermore, within-category differences are deferred until all category-level comparisons have been exhausted, rather than treated as true indifferences.

\subsection{Shifted-grid sensitivity diagnostic}
Because the ceiling convention $B_\delta(z) = \lceil z/\delta \rceil$ fixes the grid origin at zero, boundary effects may occur depending on where the grid falls. To diagnose sensitivity to arbitrary category placement, one can evaluate shifted grids $B_{\delta, a}(z) = \lceil (z-a)/\delta \rceil$ for shifts $a \in \{0, \delta/4, \delta/2, 3\delta/4\}$. If the recommended winner frequently changes under these shifts, the selection is highly sensitive to boundary placement.

\section{Theoretical Properties}
\label{sec:theory}

\begin{theorem}[Existence and completeness]
For every finite nonempty candidate set $A$, finite retained probe family and positive declared resolution vector, a ReLexTail-optimal alternative exists, and $\preceq_{\mathrm{ReLexTail}}$ is a complete preorder.
\end{theorem}
\begin{proof}
$\Psi(x)$ is a finite-dimensional vector. Lexicographic order on these vectors is a complete preorder (not a total order, as distinct alternatives can have identical $\Psi$). The pullback through $\Psi$ is a complete preorder, which has at least one minimal element on a finite non-empty set.
\end{proof}

\begin{theorem}[Positional anonymity conditional on the declared multiset]
For every permutation $\sigma$ of the probe coordinates, $\Psi(\sigma D) = \Psi(D)$.
\end{theorem}
\begin{proof}
$M$, empirical CVaR, and resolution categories depend only on symmetric quantities (order statistics). $D^\downarrow$ is permutation-invariant. Selective duplication changes the multiset and can change the result.
\end{proof}

\begin{proposition}[Strict profile monotonicity]
If $D(x) \leq D(y)$, then $x \preceq_{\mathrm{ReLexTail}} y$. If strictly less somewhere, $x \prec_{\mathrm{ReLexTail}} y$.
\end{proposition}

\begin{theorem}[Criterion-level Pareto compatibility]
If $x$ Pareto-dominates $y$, retained probes are nondecreasing, and the retained family separates criteria, then $x \prec_{\mathrm{ReLexTail}} y$.
\end{theorem}
\begin{proof}
Since $x$ Pareto-dominates $y$ and the probes are monotone, we have $D_q(x) \leq D_q(y)$ for all $q$.
By monotonicity of empirical CVaR, $T_\alpha(x) \leq T_\alpha(y)$ for all $\alpha$. The binning function $B_\delta$ is nondecreasing, so $C_\alpha(x) \leq C_\alpha(y)$. Thus weak Pareto compatibility holds. For strict Pareto compatibility, note that if the family separates criteria, $D(x) \neq D(y)$ and the exact mean $T_{100}(x) < T_{100}(y)$. Thus, even if all resolution categories tie, the strict inequality at $T_{100}$ ensures strict preference, since no earlier category can favour $y$.
\end{proof}

\begin{theorem}[Uniform replication invariance]
Let $rD$ denote the profile obtained by repeating every disappointment $r$ times. Then $D \preceq_{\mathrm{ReLexTail}} E \iff rD \preceq_{\mathrm{ReLexTail}} rE$.
\end{theorem}
\begin{proof}
The proof distinguishes the empirical distribution block from the terminal sorted-profile block. First, the empirical distribution of $rD$ is identical to $D$, meaning the empirical functionals remain unchanged: $M(rD) = M(D)$ and $T_\alpha(rD) = T_\alpha(D)$, so all resolution categories coincide. Second, if $D^\downarrow(x) <_{\mathrm{lex}} E^\downarrow(y)$, the uniformly repeated sorted vectors preserve the first differing original coordinate and its direction, preserving the lexicographic comparison.
\end{proof}

\begin{theorem}[Category stability]
Empirical upper-tail CVaR and the maximum are $1$-Lipschitz in the supremum norm. 
For any risk coordinate $R_\ell \in \{M, T_{25}, T_{50}, T_{100}\}$, define the distance to the nearest category boundary as $\mu_\ell(x) = \operatorname{dist}(R_\ell(x), \delta_\ell \mathbb{Z})$. Let $\mu_{\mathrm{cell}} = \min_{x\in A, \ell} \mu_\ell(x)$. If a perturbation satisfies $\eta < \mu_{\mathrm{cell}}$, then every candidate remains in the same category on every category coordinate.
\end{theorem}
\begin{proof}
For the maximum $M$, the $1$-Lipschitz property is immediate since $|M(D) - M(\widehat{D})| \leq \max_q |D_q - \widehat{D}_q| \leq \eta$. For empirical upper-tail CVaR, by definition of the perturbation bound, $D - \eta\mathbf{1} \leq \widehat{D} \leq D + \eta\mathbf{1}$. By monotonicity and translation equivariance of CVaR, $T_\alpha(D) - \eta \leq T_\alpha(\widehat{D}) \leq T_\alpha(D) + \eta$, yielding $|T_\alpha(D) - T_\alpha(\widehat{D})| \leq \eta$. Because the shift in every $R_\ell$ is at most $\eta < \mu_{\mathrm{cell}}$, no value crosses a category boundary $\delta_\ell \mathbb{Z}$, preserving all categories.
\end{proof}

\begin{theorem}[Winner stability]
Once categories are fixed, comparison proceeds through exact coordinates. 
If every rival $y$ is separated from the winner $x^\star$ at the first exact coordinate $j(y)$ immediately following the category block, then no preceding exact tie exists. Let $\Delta_{\mathrm{exact}} = \min_{y \neq x^\star} [\Psi_{j(y)}(y) - \Psi_{j(y)}(x^\star)]$. If $2\eta < \Delta_{\mathrm{exact}}$ and all categories remain unchanged, the winner remains unchanged.
\end{theorem}
\begin{proof}
Since separation occurs at the first exact coordinate, the winner and rival are tied on all categories, and strictly separated at $j(y)$. The gap between them is at least $\Delta_{\mathrm{exact}}$. A perturbation of at most $\eta$ in the underlying profile shifts each exact coordinate by at most $\eta$, changing the gap by at most $2\eta$. Since $2\eta < \Delta_{\mathrm{exact}}$, the strict inequality at $j(y)$ is preserved, maintaining $x^\star$ as the winner.
\end{proof}
